% =============================================================================
% 5. ML ENHANCEMENT COMPONENT
% Target: 1 page (~500 words)
% Source: cap4 §4.6 (sec:ml-component)
% Status: SKELETON
% =============================================================================

\section{Machine-Learning Enhancement Component}
\label{sec:ml-component}

The ML component $\mathcal{S}_{\text{ML}}$ runs in parallel to the
sequential pipeline and produces a case-level predictor $\hat{f}$. Its
contract is summarized in Table~\ref{tab:contracts}; its novelty is the
\emph{family of features} it consumes, half of which are PM- and
Markov-derived and are not available to traditional MSR baselines.

% ---------- §5.1 Feature engineering ----------
\subsection{Feature Engineering}
\label{sec:ml-features}

\begin{table}[ht]
\centering
\caption{Feature definitions. Features 1--8 are derived from the PM and
Markov stages; features 9--15 are traditional MSR features used as
baseline.}
\label{tab:features}
\small
\begin{tabular}{@{}clll@{}}
\toprule
\textbf{\#} & \textbf{Feature} & \textbf{Source} & \textbf{Definition} \\
\midrule
\multicolumn{4}{@{}l}{\textit{Process-derived (novel)}} \\
1 & Current state          & $\mathcal{L}$       & $\mu(e_{\text{last}})$ \\
2 & Sojourn in current state & $\mathcal{L}$     & $t_{\text{obs}} - \tau(e_{\text{last}})$ \\
3 & Transition count       & $\mathcal{L}$       & $|\sigma_c^{\text{partial}}|$ \\
4 & Rework count           & $\mathcal{L}$       & regression count \\
5 & Conformance score      & $M,\mathcal{L}$     & $\phi_{\text{conf}}(c)$ \\
6 & Markov expected remaining & $N$              & $t_{\mu(e_{\text{last}})}$ \\
7 & Completion probability & $B$                 & $B_{\mu(e_{\text{last}}),\texttt{done}}$ \\
8 & Sojourn ratio          & $\hat{F}$           & sojourn/$\bar{d}_{\mu(e_{\text{last}})}$ \\
\midrule
\multicolumn{4}{@{}l}{\textit{Traditional MSR (baseline)}} \\
9--15 & Files / lines / commits / authors / priority / type / etc.\
     & Git, Jira & --- \\
\bottomrule
\end{tabular}
\end{table}

Three feature sets are evaluated in the companion empirical paper:
\textbf{PM-only} (1--8), \textbf{MSR-only} (9--15), and
\textbf{Combined} (1--15).

% ---------- §5.2 Label construction & training protocol ----------
\subsection{Label Construction and Training Protocol}
\label{sec:ml-training}

\begin{definition}[Outcome Label]
\label{def:label}
$y_c = 1$ if $T_c > \tau_{\text{threshold}}$
(\texttt{delayed}); $y_c = 0$ otherwise (\texttt{on\_time}).
\end{definition}

Three threshold strategies are supported:
\emph{percentile-based} ($\tau_{\text{threshold}} = \hat{Q}_{0.75}$),
\emph{SLA-based}, and \emph{relative}
($2 \cdot \tilde{d}_{\text{type}}$).

The training protocol enforces a strict temporal hold-out:
$\mathcal{L}_{\text{train}}$ contains all cases completed before
$t_{\text{cut}}$, $\mathcal{L}_{\text{test}}$ contains only cases that
\emph{start} after $t_{\text{cut}}$. Candidate models
$\{\text{RF}, \text{XGBoost}, \text{LogReg}\}$ are compared by
temporal cross-validation with expanding
windows~\cite{tashman2000out, hyndman2021forecasting}; the best model
$\hat{f}$ is evaluated on $\mathcal{L}_{\text{test}}$ via AUC-ROC,
F1, Brier score, and the feature-importance vector $\mathbf{w}$ is
reported.

% ---------- §5.3 Integration with Monte Carlo ----------
\subsection{Integration with Monte Carlo Forecasts}
\label{sec:ml-integration}

\begin{definition}[Deviation Score]
\label{def:deviation}
\begin{equation}
\label{eq:deviation}
\delta_c = \hat{f}(\mathbf{x}_c) -
           \hat{P}_{\text{MC}}(\text{delayed} \mid s_0 = s).
\end{equation}
\end{definition}

Monte Carlo (§\ref{sec:stages-s4}) produces population-level
distributional forecasts; the ML component produces case-level point
forecasts; $\delta_c$ identifies cases departing from the population
baseline.
